% ==============================================================================
% T0/FFGFT Theory: Document 240
% Title: AI Detectors and Ad Hominem Argumentation
% Subtitle: Structural Kinship of Old and New Manipulation Patterns
% Author: Johann Pascher
% Date: May 2026
% Reference: methodologically self-contained, loosely related to Doc 207 §10.3
% Occasion: IPI mailing list, discussion with Peter Austin (May 2026)
% ==============================================================================

\section*{Preamble}

This text grew out of a discussion on the mailing list of the
\emph{Information Physics Institute} (May 2026). Peter Austin
reported that his own emails are regularly marked by AI detectors as
more than 50\,\% AI-generated -- even when written in personal
correspondence and without any AI assistance. The initial question
in the group was practical: what to do about these false positives.

On closer examination, the question turns out to be not primarily
technical but methodological and psychological. AI detectors are
not the beginning of a new phenomenon. They stand at the provisional
end of an old line: the argumentation patterns that reject without
engaging with the substance. This document describes the structural
kinship and names the psychological mechanism that drives the pattern.

% ==============================================================================
\section{The contradiction in the demand}
% ==============================================================================

At the centre of the current debate stands an open contradiction:

\begin{itemize}[leftmargin=2em,itemsep=2pt,topsep=2pt]
  \item On one side we demand precise, error-free, well-structured
        texts. This has always been the standard of scientific and
        professional communication.
  \item On the other side, \emph{the absence of errors} is now
        itself suspicious. A text with low perplexity, uniform
        sentence structure, and clean transitions is flagged by
        detectors as AI.
\end{itemize}

Taken literally, this is an \textbf{instruction to write worse}:
produce enough roughness to pass through, enough variation to avoid
being read as a machine. This cannot be a methodological standard.

The evidence for this skew is well documented. The Stanford study
by Weixin Liang, Mert Yuksekgonul, Yining Mao, Eric Wu and James Zou
-- \emph{GPT detectors are biased against non-native English writers}
(Patterns, Vol.~4, No.~7, Article 100779, 14 July 2023;
arXiv:2304.02819) -- showed: seven widely used GPT detectors
classified \emph{more than half} of TOEFL essays by non-native
English learners as \enquote{AI-generated}, while the same detectors
achieved near-perfect accuracy on US 8th-grade essays. The reason:
non-native texts typically have lower perplexity, because vocabulary
is more limited and sentence structures simpler -- exactly the
feature they share with machine-generated text.

OpenAI itself -- the maker of ChatGPT -- published in January 2023
an \emph{AI Classifier} that was supposed to solve precisely this
problem. Six months later, on 20 July 2023, the tool was withdrawn.
The official reason given by OpenAI: only \emph{26\,\%} of
genuinely AI-generated texts were correctly identified, and in
\emph{9\,\%} of cases human texts were falsely flagged as AI. With
a hit rate below 50\,\%, a detector is worse than a coin flip.

Anyone who writes carefully as a technician, scientist, or
non-native speaker necessarily produces clearer and more structured
text than a casual writer. It is precisely this writing that is
systematically flagged.

% ==============================================================================
\section{Substance, not style}
% ==============================================================================

The real distinction between human and machine-generated content
does not lie in style, but in substance.

\textbf{AI texts are characterised not by their form, but by a
specific substantive weakness:} they are often coherent-sounding
arguments built on \emph{hallucinated facts}, \emph{fabricated
citations}, and \emph{confidently presented but vacuous derivations}.
The linguistic polish is real; the substance is invented.

A critical reader who understands the subject recognises this
weakness -- not from surface markers, but because the claims fail
on closer inspection. A statistical machine cannot in principle see
this difference, because it is not a surface feature.

\textbf{Consequence:} whoever wants to judge content must be able
to understand content. Detectors cannot do this. They can only
count symptoms.

% ==============================================================================
\section{The detector user as the actual misuser}
% ==============================================================================

Here an inversion appears that is rarely articulated in the debate.
Whoever uses an AI detector instead of reading the text themselves
\textbf{delegates their own judgement to an algorithm}.

The detector user is thereby more deeply dependent on AI than the
author who consciously uses AI as a tool:

\begin{center}
\renewcommand{\arraystretch}{1.4}
{\small
\begin{tabular}{>{\raggedright\arraybackslash}p{3.4cm}>{\raggedright\arraybackslash}p{3.6cm}>{\raggedright\arraybackslash}p{3.4cm}}
\toprule
\textbf{Role} & \textbf{What the person does} & \textbf{Who holds the judgement} \\
\midrule
Author with AI tool & substance their own, language with AI help & human \\[3pt]
Reviewer with detector & reads score, takes judgement from machine & machine \\
\bottomrule
\end{tabular}
}
\renewcommand{\arraystretch}{1.0}
\end{center}

The reproach \enquote{you used AI} therefore often comes from
exactly the person who has already handed over their own critical
function to AI. They accuse the author of what they themselves do
-- without being aware that they are doing it.

% ==============================================================================
\section{AI as filter, not as oracle}
% ==============================================================================

An important distinction is needed. Even careful knowledge workers
use AI -- for summaries, for language formalisation, for pre-sorting
large amounts of text. This is not the problem. The problem is
\emph{the kind} of use.

\textbf{AI as filter (triage):}
\begin{itemize}[leftmargin=2em,itemsep=2pt,topsep=2pt]
  \item AI reads large amounts of text and summarises.
  \item The human reads the summary and decides: does this match,
        or do I need to look more closely?
  \item In case of ambiguity or unusual findings: the human goes
        into the detail.
  \item The final judgement always rests with the human.
\end{itemize}

\textbf{AI as oracle (delegation):}
\begin{itemize}[leftmargin=2em,itemsep=2pt,topsep=2pt]
  \item AI delivers a score, a classification, a truth.
  \item The score is accepted.
  \item No own examination of substance.
  \item No willingness to read for oneself in case of irregularities.
\end{itemize}

AI use as \textbf{filter} is methodologically unproblematic -- it
corresponds to the long-established practice of having larger
amounts of text pre-sorted by abstracts, bibliographies, or
assistants. AI use as \textbf{oracle} is methodologically fatal:
it replaces judgement instead of supporting it.

Detectors are the paradigm case of oracle use. They deliver a
number that can be accepted or ignored -- but no substance that
could be examined.

% ==============================================================================
\section{Structural impossibility of substantive judgement}
% ==============================================================================

Here lies the technical point that pulls the rug from under the
entire detector construction.

An AI detector implicitly claims: \enquote{I can recognise whether
a text comes from an AI.} For this it would have to do something
that AI itself cannot do:

\begin{itemize}[leftmargin=2em,itemsep=2pt,topsep=2pt]
  \item AI does not know whether a claim is true.
  \item AI does not know whether an argument holds.
  \item AI does not recognise its own hallucinations.
\end{itemize}

\textbf{If AI itself is overwhelmed at the substance level, a
detector -- which is also AI -- cannot judge substance either.}
It can only measure surface features: word probabilities
(perplexity), sentence-length variation (burstiness), frequency
of certain transition words.

The construction dissolves itself logically:

\begin{quote}
Whoever wants to judge substance must be able to understand it.
Detectors are AI. AI cannot understand substance. So detectors
cannot judge what they claim to judge.
\end{quote}

A detector that could really judge substance would no longer be an
AI -- it would be a \emph{critically thinking human with subject
expertise}. But that is precisely what the detector is supposed to
replace. The claim is self-contradictory.

\subsection{An analogy}

No one would believe someone who said: \enquote{I can tell whether
a maths book is written by an incompetent author -- without knowing
any mathematics myself. I just measure the frequency of certain
words.} But this is exactly what AI detectors do: they claim to
judge \textbf{substance} through \textbf{form statistics}.

This is the same category error as in phrenology (skull shape
equals character) or polygraphs (sweat production equals truth):
\emph{measuring surface correlates and presenting it as measurement
of the underlying quantity.}

% ==============================================================================
\section{Kinship with ad hominem argumentation}
% ==============================================================================

Up to here the argument has been technical. Now comes the structural
point: AI detectors are not a new phenomenon.

\textbf{They do to texts what ad hominem argumentation does to
people.}

\begin{center}
\renewcommand{\arraystretch}{1.3}
{\small
\begin{tabular}{>{\raggedright\arraybackslash}p{5.0cm}>{\raggedright\arraybackslash}p{5.6cm}}
\toprule
\textbf{Ad hominem argumentation} & \textbf{AI detector} \\
\midrule
attacks the speaker, not the argument & rates the form, not the content \\[2pt]
who said it? & who (or what) wrote it? \\[2pt]
bypasses substantive engagement & bypasses substantive examination \\[2pt]
convenient -- no understanding needed & convenient -- no reading needed \\[2pt]
replaces thinking with suspicion & replaces reading with a score \\
\bottomrule
\end{tabular}
}
\renewcommand{\arraystretch}{1.0}
\end{center}

Both follow the same pattern: \textbf{shifting evaluation from the
matter itself to a property of the source.} Instead of testing
whether something is right, the question becomes where it comes
from -- and from origin a conclusion about validity is drawn.

In logic this has been known for centuries as the \emph{genetic
fallacy}. Aristotle described ad hominem arguments in his
\emph{Sophistical Refutations} as fallacies. Arthur Schopenhauer
addresses the related \emph{argumentum ad personam} in his
\emph{Eristic Dialectic} as a pseudo-argument that bypasses the
matter. The mechanism is 2400 years old; the technical completion
in software form is new.

% ==============================================================================
\section{Cognitive dissonance as the driver of the pattern}
% ==============================================================================

Why does this pattern work so reliably? The answer comes from Leon
Festinger's \emph{A Theory of Cognitive Dissonance} (Stanford
University Press, 1957), one of the most influential works of
20th-century social psychology.

\subsection{Festinger's basic observation}

Festinger developed the theory from an unusual empirical
observation. In \emph{When Prophecy Fails} (with Henry W. Riecken
and Stanley Schachter, University of Minnesota Press, 1956) he
documents the behaviour of a doomsday cult that had predicted a
concrete date for the end of the world. When the date passed
without anything happening, the members' belief \emph{strengthened}
rather than collapsing. They interpreted the failure of the
prophecy as proof that their faith community had saved the world.

Festinger's conclusion: when two convictions are incompatible,
psychological pressure (\emph{dissonance}) arises. This pressure
does not necessarily lead to correction of the false belief --
often it leads to \emph{reinterpretation of reality} so that the
original position remains confirmed.

\subsection{Application to the detector context}

In the AI-detector situation, the following dissonance typically
arises in the reviewer:

\begin{center}
\renewcommand{\arraystretch}{1.3}
{\small
\begin{tabular}{>{\raggedright\arraybackslash}p{5.0cm}>{\raggedright\arraybackslash}p{5.6cm}}
\toprule
\textbf{Conviction A} & \textbf{Conviction B (triggering)} \\
\midrule
\enquote{I am a competent reviewer} & \enquote{Before me lies a text clearer than what I would write} \\[2pt]
\enquote{I belong to those who think critically} & \enquote{Someone here is apparently thinking more clearly than I am} \\[2pt]
\enquote{AI makes people dumber} & \enquote{But this text does not seem dumb -- it seems cleverer than usual} \\
\bottomrule
\end{tabular}
}
\renewcommand{\arraystretch}{1.0}
\end{center}

The dissonance is uncomfortable. There are two ways to resolve it:

\textbf{Path~1: Substantive engagement} (strenuous).\\
Reconsider one's own position; acknowledge that the other has done
something good; adjust one's own standards; revise one's own
position if necessary.

\textbf{Path~2: Source devaluation} (convenient).\\
Mark the text as inauthentic -- \enquote{this is not really from
them, this is AI}. Self-perception remains intact. No own revision
needed.

The AI detector provides the \emph{technical justification} for
Path~2. It allows the uncomfortable text to be disqualified without
having to engage with it. The score looks scientific, objective,
neutral -- and provides exactly what the reviewer needs to resolve
the dissonance without questioning themselves.

\subsection{Three functions in one}

Dissonance reduction through source devaluation fulfils three
functions simultaneously:

\begin{itemize}[leftmargin=2em,itemsep=2pt,topsep=2pt]
  \item \textbf{Self-esteem protection:} \enquote{I am not less
        competent -- they just cheated.}
  \item \textbf{Effort avoidance:} \enquote{I do not have to engage
        with this.}
  \item \textbf{Social positioning:} \enquote{I am one of the
        critical ones who spot AI cheating.}
\end{itemize}

All three functions are \emph{cheap} to obtain -- a single click on
a detector. The honest alternative -- substantive examination -- is
\emph{expensive}: time, concentration, the risk of having to revise
one's own position.

\subsection{Amplification through pseudo-objectivity}

Classical ad hominem arguments can be directly rebutted by the
addressee -- one can reply, contextualise, defend oneself. AI
detectors have three additional features that make dissonance
reduction even more effective:

\begin{enumerate}[leftmargin=2em,itemsep=2pt,topsep=2pt]
  \item \textbf{Pseudo-objectivity:} The number looks scientific,
        not personal. The reviewer can say: \enquote{I am not
        attacking you -- the algorithm says this.}
  \item \textbf{Responsibility offloading:} Nobody feels personally
        responsible for the verdict. \enquote{The computer decided.}
  \item \textbf{Irrefutability:} The addressee \emph{cannot prove}
        that they did not use AI. Negation is not provable.
\end{enumerate}

These three features make the AI detector an \emph{ideal tool for
dissonance reduction}: maximum effect with minimal personal
investment and minimal responsibility.

\subsection{A bitter point}

From this follows a paradoxical consequence: the \textbf{clearer
and more competent} a text is, the \textbf{greater the dissonance}
for a less self-assured reader -- and the \textbf{stronger the
motive} to mark the text as AI. Detectors thus punish exactly the
competence they claim to protect. Not by accident, but by
psychological function.

% ==============================================================================
\section{The whole family of rhetorical manipulation patterns}
% ==============================================================================

The kinship is not confined to the ad hominem fallacy. AI detectors
in fact reproduce a \emph{whole family} of classical manipulation
techniques -- each in technical form.

\subsection{Tone policing}
\textbf{Human:} \enquote{You are getting emotional -- one cannot
take this seriously.}\\
\textbf{Detector:} \enquote{You write too smoothly -- one cannot
take this seriously.}

Style becomes the disqualifier, regardless of content.

\subsection{Strawman argument}
The \emph{human} distorts the opponent's position into a simpler,
easier-to-refute caricature. The \emph{detector} reduces the text
to statistical surface features that have nothing to do with the
actual content. In both cases a simplified caricature is judged,
not the original.

\subsection{Burden-of-proof reversal}
\textbf{Human:} \enquote{Prove to me that you thought this up
yourself.}\\
\textbf{Detector:} \enquote{Prove to me that you didn't use AI.}

In both cases the opposite is taken as default. The addressee has
to prove their innocence -- which is structurally impossible
because \emph{non-existence is not provable}.

\subsection{Argumentum ad verecundiam}
\textbf{Human:} \enquote{Professor~X said so, therefore it must be
true.}\\
\textbf{Detector:} \enquote{The algorithm says 87\,\%, therefore it
must be true.}

An authority is placed above the content. With AI detectors the
authority is a machine -- which sharpens the inversion.

\subsection{Argumentum ad populum}
\textbf{Human:} \enquote{Everyone knows that\ldots}\\
\textbf{Detector:} \enquote{This looks like typical AI style, all
the models recognise the pattern.}

Consensus (or learned pattern) replaces argument.

\subsection{Gaslighting}
\textbf{Human:} \enquote{You overestimate yourself -- this cannot
come from you.}\\
\textbf{Detector:} \enquote{The score shows: this cannot come from
you.}

The addressee is denied their own authorship. They are made to
doubt themselves. This is the most destructive form -- it targets
not the claim, but the person.

% ==============================================================================
\section{The historical line}
% ==============================================================================

The patterns described here are not new. They are documented in
every epoch:

\begin{itemize}[leftmargin=2em,itemsep=2pt,topsep=2pt]
  \item \textbf{Ancient rhetoric:} Aristotle described ad hominem
        arguments in the \emph{Sophistical Refutations} as
        fallacies.
  \item \textbf{Middle Ages:} \emph{Argumentum ad auctoritatem}
        (appeal to authorities instead of own examination) as a
        scholastic argumentation form.
  \item \textbf{Inquisition:} suspicion logic -- whoever thinks
        like this is suspect.
  \item \textbf{20th-century totalitarianism:} \enquote{whoever
        speaks like this is an enemy.}
  \item \textbf{Present:} AI detectors -- \enquote{whoever writes
        like this is not real.}
\end{itemize}

It is always the same pattern: when substantive engagement becomes
difficult, one falls back on source criticism. The \textbf{historical
constancy} is the most important observation of this document -- it
shows that the real problem is not technological but
\emph{anthropological}.

% ==============================================================================
\section{What would actually help}
% ==============================================================================

If detectors do not help but worsen the problem, what does?

\textbf{Substantive reading.} A competent reader who understands
the subject recognises hallucinations because they understand the
substance. No score required. What is required is willingness to
engage with the content.

\textbf{Reproducible verification.} Where content is checkable, it
can be checked -- through code, data, scripts that one can run
oneself. This is not a detector function; this is a
\emph{content function}.

\textbf{Provenance documentation.} Drafts, voice recordings,
revision history. Anyone who makes their working process
transparent shifts the question from \enquote{prove you didn't use
AI} to \enquote{here is the workflow, see for yourself}. The second
question can be answered.

\textbf{Open discipline.} Use AI consciously as a tool, not as an
oracle. Mark your own position as your own. Acknowledge AI
assistance openly rather than hiding it.

None of these points needs a detector. All four need \emph{competent,
honest people who are willing to do the work of understanding}.

% ==============================================================================
\section{Diagnosis}
% ==============================================================================

The AI-detector debate is not primarily technical. It is the
visible peak of a cultural tendency that replaces the strenuous
business of understanding with the convenient practice of source
criticism.

In one sentence:

\begin{quote}
\textbf{Ad hominem argumentation is the AI detector for people. The
AI detector is ad hominem argumentation for texts. Both are the
refusal of the matter, dressed up as concern for the source.}
\end{quote}

The psychological driver -- cognitive dissonance and its resolution
through source devaluation -- has been well documented since
Festinger 1956/1957. The historical constancy of the pattern
reaches back to Aristotle. The technical completion in software
form is new; the underlying problem is old.

Whoever trusts an AI checker more than their own critical judgement
has ceased to be a reviewer -- and has themselves become an
appendage of a machine. That is the real misuser of AI.

% ==============================================================================
\section{Sources}
% ==============================================================================

\begin{itemize}[leftmargin=2em,itemsep=2pt,topsep=2pt]
  \item Liang, W., Yuksekgonul, M., Mao, Y., Wu, E., Zou, J.
        (2023). \emph{GPT detectors are biased against non-native
        English writers}. Patterns Vol.~4, No.~7, Article 100779.
        DOI:~10.1016/j.patter.2023.100779; arXiv:2304.02819.
  \item OpenAI (2023). \emph{New AI classifier for indicating
        AI-written text}. Blog post of 31~January 2023, updated on
        20~July 2023 with the note announcing the withdrawal of the
        tool due to low accuracy.
  \item Festinger, L. (1957). \emph{A Theory of Cognitive
        Dissonance}. Stanford University Press.
  \item Festinger, L., Riecken, H.~W., Schachter, S. (1956).
        \emph{When Prophecy Fails: A Social and Psychological
        Study of a Modern Group that Predicted the Destruction of
        the World}. University of Minnesota Press.
  \item Aristotle. \emph{Sophistical Refutations} (Sophistikoi
        Elenchoi). Classical account of rhetorical fallacies.
  \item Schopenhauer, A. \emph{The Art of Being Right / Eristic
        Dialectic} (published posthumously). Treatment of the
        \emph{argumentum ad personam}.
\end{itemize}

\bigskip

\textbf{Note on the genesis of this document.} The text was written
in May 2026, prompted by a discussion on the mailing list of the
\emph{Information Physics Institute} with Peter Austin. The argument
is my own thinking. Language formalisation was done with AI
assistance -- in exactly the manner described in §5 as \emph{AI as
filter, not as oracle}. An open discipline: substance from the
author, language with tool support, judgement with the human.
